
\titulo{5}

Calamar de la noche,\\
despiadada marítima criatura\\
que sumerge nuestras embarcaciones,\\
señor de los naufragios\\
y de enormes ojos desorbitados:\\
invoco tu presencia con temblor en los labios.\\
En mi boca vive sólo tu nombre,\\
tu cara puebla todos mis horrores,\\
tu olor es el perfume del palosanto.\\
Tus prénsiles tentáculos\\
amenazan la vaga luz del alba.\\
Tu fosa ha sido abierta,\\
las lágrimas que plañes han salado los mares,\\
tu oscuridad relumbra\\
fosforescente en las profundidades\\
con la luminiscencia de los ángeles\\
entonando cánticos ancestrales.\\
Calamar de la noche:\\
las laboriosas civilizaciones\\
resecas ya por el natrón del tiempo\\
veneraron tu náutica presencia\\
en ánforas e intrincados mosaicos.\\
Calamar de la noche,\\
señor de los naufragios,\\
bajándote la luna\\
encomiendo mi navío a tus manos:\\
traigas la noche al día,\\
ensombrezcas nuestros diarios caminos,\\
nos protejan de los vientos tus brazos,\\
los miedos borre el aura de tu llanto.\\

